% =============================================================
% บทที่ 1 - บทนำ
% หัวข้อ: การวิเคราะห์และทำนายผลการแข่งขันฟุตบอลลีกยุโรปด้วยกระบวนการวิทยาการข้อมูล
%        เพื่อการประเมินมูลค่าการเดิมพัน
% =============================================================

% -----------------------------------------------------------
\section{ความเป็นมาและความสำคัญของปัญหา}
% -----------------------------------------------------------

ปัจจุบัน ตลาดการเดิมพันกีฬาฟุตบอลถือเป็นตลาดการเงินทางเลือกที่มีมูลค่ามหาศาลทั่วโลก โดยเจ้ามือรับพนัน (Bookmakers) ได้ประมวลผลสถิติเชิงลึกร่วมกับแบบจำลองเชิงสถิติขั้นสูงเพื่อกำหนดราคาอัตราต่อรอง (Odds) ซึ่งสะท้อนถึงความน่าจะเป็นของผลการแข่งขัน อย่างไรก็ตาม ราคาต่อรองดังกล่าวมักถูกบวกส่วนต่างกำไรของเจ้ามือ (Overround / Vig) และมีการเปลี่ยนแปลงตามกระแสการเดิมพันของตลาด ส่งผลให้นักเดิมพันส่วนใหญ่ที่ตัดสินใจโดยใช้อารมณ์หรือสถิติพื้นฐานไม่สามารถสร้างผลตอบแทนที่ยั่งยืนได้ในระยะยาว 

การพัฒนาแบบจำลองการทำนายผลการแข่งขันฟุตบอลยุโรป (ฤดูกาล 2014--2015 ถึง 2025--2026) ด้วยกระบวนการวิทยาการข้อมูล (Data Science Process) และการเรียนรู้ของเครื่อง (Machine Learning) จึงเข้ามามีบทบาทสำคัญในการประเมิน \textbf{ความน่าจะเป็นที่แท้จริง (True Probabilities)} ของผลการแข่งขัน 3 ทางเลือก ได้แก่ เจ้าบ้านชนะ (Home Win: H), เสมอ (Draw: D) และทีมเยือนชนะ (Away Win: A) เพื่อนำไปเปรียบเทียบกับความน่าจะเป็นที่แฝงอยู่ในราคาต่อรองของตลาด (Implied Probabilities) สำหรับค้นหาโอกาสการเดิมพันที่มีมูลค่าคาดหมายเป็นบวก (Value Betting / Positive Expected Value: +EV)

อย่างไรก็ตาม การสร้างแบบจำลองเพื่อประเมินศักยภาพการแข่งขันให้เหนือกว่าราคาตลาดมีความท้าทายสูงมาก เนื่องจากงานวิจัยในอดีต \cite{hubacek2021forty} ชี้ว่าแบบจำลองที่พึ่งพาเพียงสถิติผลย้อนหลังมักมีความแม่นยำจำกัด โครงงานนี้จึงมุ่งแก้ปัญหาดังกล่าวด้วยการบูรณาการ \textbf{ข้อมูลบริบทเชิงลึกก่อนการแข่งขัน (Pre-match Contextual Features)} ครอบคลุม 5 ลีกใหญ่ของยุโรป ได้แก่ คะแนนความแข็งแกร่ง Elo Rating, ดัชนีความเหนื่อยล้าสะสมของผู้เล่นตัวจริง, ความหนาแน่นของตารางแข่งขัน (รวมถึงฟุตบอลถ้วยและรายการยุโรป) และฟอร์มรายบุคคลของผู้เล่น 11 ตัวจริง (เช่น xG และ xA ต่อ 90 นาที) เพื่อนำมาสร้างแบบจำลองการทำนาย แล้วนำความน่าจะเป็นที่ได้ไปเข้ากลยุทธ์การจัดสรรเงินทุน (Staking Strategies) เช่น Kelly Criterion และ Fixed Staking พร้อมทดสอบย้อนหลัง (Backtesting) เพื่อวัดผลตอบแทนการลงทุน (ROI) ในสถานการณ์ที่สมจริง

% -----------------------------------------------------------
\section{วัตถุประสงค์}
% -----------------------------------------------------------

\begin{enumerate}[label=\arabic*)]
    \item เพื่อศึกษาและพัฒนากระบวนการวิทยาการข้อมูลอย่างครบวงจร ตั้งแต่การรวบรวมข้อมูลหลายแหล่ง (สถิติการแข่งขัน, แผนการเล่น และราคาต่อรอง) การทำความสะอาดข้อมูล ไปจนถึงการพัฒนาแบบจำลองทำนายผลลัพธ์การแข่งขัน
    \item เพื่อสร้างและวิเคราะห์ตัวแปรก่อนการแข่งขัน (Pre-match Features) ได้แก่ คะแนน Elo, ดัชนีความเหนื่อยล้า, ความหนาแน่นของตารางแข่งขัน และฟอร์มผู้เล่นตัวจริง 11 คนแรก ที่มีผลต่อโอกาสชนะและความคุ้มค่าในการเดิมพัน
    \item เพื่อพัฒนาระบบคำนวณและเปรียบเทียบความน่าจะเป็นของแบบจำลองกับราคาอัตราต่อรองของตลาด (Market Odds Analysis) เพื่อค้นหาโอกาสการเดิมพันที่มีมูลค่า (Value Betting)
    \item เพื่อประเมินประสิทธิภาพของแบบจำลองและประเมินผลกำไรการลงทุนผ่านระบบทดสอบย้อนหลัง (Backtesting) ด้วยกลยุทธ์การบริหารเงินทุน (Staking Strategies) รูปแบบต่าง ๆ
\end{enumerate}

% -----------------------------------------------------------
\section{ขอบเขตและข้อจำกัดของโครงงาน}
% -----------------------------------------------------------

\subsection{ขอบเขตของโครงงาน}

\begin{enumerate}[label=\arabic*)]
    \item \textbf{ขอบเขตด้านข้อมูลและช่วงเวลา:} ใช้ชุดข้อมูลการแข่งขันฟุตบอลย้อนหลังครอบคลุมฤดูกาล 2014--2015 ถึง 2025--2026 และฤดูกาล 2026--2027 ที่กำลังแข่งขันอยู่ (ถึงวันที่ 11 ตุลาคม 2026) จากแหล่งข้อมูลสาธารณะ ได้แก่
    \begin{itemize}[labelindent=0pt]
        \item \textit{Understat}: รวบรวมสถิติผลการแข่งขัน ค่าความคาดหวังในการทำประตู (xG, xA) และสถิติเชิงลึกรายบุคคลของผู้เล่น \cite{understat2024}
        \item \textit{Transfermarkt}: รวบรวมปฏิทินการแข่งขันครบทุกรายการ (รวมฟุตบอลถ้วยและรายการยุโรป) และแผนการเล่น (Formation) ของผู้เล่น 11 ตัวจริง \cite{transfermarkt2024}
        \item \textit{Football-Data.co.uk}: รวบรวมข้อมูลผลการแข่งขัน ยิงประตู และราคาอัตราต่อรองย้อนหลัง (Historical Betting Odds Data Service) จากเจ้ามือรับพนันหลายแห่ง สำหรับใช้ทดสอบแบบจำลองและประเมินผลตอบแทนทางการเงิน
    \end{itemize}
    ยกเว้นการทดสอบสมมติฐานทางสถิติ (ความได้เปรียบของเจ้าบ้านและความสัมพันธ์ของคะแนน Elo กับผลต่างประตู) ซึ่งใช้ผลการแข่งขันพรีเมียร์ลีกอังกฤษฤดูกาล 2000--2001 ถึง 2024--2025 เพื่อให้มีจำนวนนัดมากพอ

    \item \textbf{ขอบเขตด้านลีกการแข่งขัน:} ระบบการรวบรวมข้อมูลและกระบวนการประมวลผลตัวแปร (Feature Pipeline) ครอบคลุม 5 ลีกใหญ่ยุโรป ได้แก่ พรีเมียร์ลีก อังกฤษ, ลาลีกา สเปน, บุนเดสลีกา เยอรมนี, กัลโช่ เซเรียอา อิตาลี และลีกเอิง ฝรั่งเศส

    \item \textbf{ขอบเขตด้านกระบวนการและระบบโครงสร้าง:} พัฒนาท่อส่งข้อมูลแบบโมดูลาร์ด้วยภาษา Python ประกอบด้วย กระบวนการจัดเก็บข้อมูล, การทำความสะอาดและสร้างตัวแปรพร้อมการตรวจสอบสัญญาโครงสร้างข้อมูล (Pandera Data Contracts), การวิเคราะห์การเดิมพันและการทดสอบย้อนหลัง (Backtesting), การวิเคราะห์ราคาตลาด (Market \& Odds Analysis), ระบบทดสอบคำสั่งอัตโนมัติ และระบบเชื่อมต่อบริการทำนายผล (FastAPI)

    \item \textbf{ขอบเขตด้านตัวแปรและการวิเคราะห์ตลาด (Market Analysis):}
    \begin{itemize}[labelindent=0pt]
        \item ใช้เฉพาะตัวแปรที่รับรู้ได้ล่วงหน้าก่อนการแข่งขัน (Pre-match Data) ได้แก่ คะแนน Elo, ดัชนีความเหนื่อยล้า, ความหนาแน่นตารางแข่งขัน และฟอร์มผู้เล่น 11 ตัวจริง
        \item วิเคราะห์ราคาต่อรองอัตราจ่าย (Odds) เพื่อคำนวณค่าส่วนต่างเจ้ามือ (Overround / Vig) และแปลงเป็นความน่าจะเป็นของตลาด (Implied Probabilities)
    \end{itemize}

    \item \textbf{ขอบเขตด้านแบบจำลองและการประเมินผลกำไร:}
    \begin{itemize}[labelindent=0pt]
        \item เปรียบเทียบแบบจำลอง Multinomial Logistic Regression, Gradient-boosted Trees (HistGradientBoosting) และแบบจำลองจำนวนประตู Poisson ในการทำนายความน่าจะเป็น และเปรียบเทียบกับความน่าจะเป็นจากราคาของ Pinnacle
        \item ทดสอบชุดกลยุทธ์การเดิมพันที่ลงทะเบียนไว้ล่วงหน้า พร้อมการปรับแก้การทดสอบหลายครั้ง และแผนการวางเงิน (Staking Plans) ได้แก่ การวางเงินคงที่ (Fixed Staking), Kelly Criterion / Fractional Kelly และ Martingale
        \item แบ่งชุดข้อมูลทดสอบตามลำดับเวลา (Temporal Split / Walk-forward) เพื่อป้องกันปัญหาข้อมูลรั่วไหลจากอนาคต (Data Leakage)
    \end{itemize}

    \item \textbf{ตัวชี้วัดประสิทธิภาพ:} 
    \begin{itemize}[labelindent=0pt]
        \item \textit{ด้านแบบจำลอง:} Log Loss, Ranked Probability Score (RPS), Accuracy และ Permutation Importance
        \item \textit{ด้านการลงทุน/เดิมพัน:} อัตราผลตอบแทนจากการลงทุน (Yield / ROI), ผลกำไรสะสม (Net Profit/Loss) และระดับความเสี่ยงการลดลงของเงินทุน (Drawdown)
    \end{itemize}
\end{enumerate}

\subsection{ข้อจำกัดของโครงงาน}

\begin{enumerate}[label=\arabic*)]
    \item ราคาต่อรองที่ใช้ในการทดสอบย้อนหลังอ้างอิงจากข้อมูล Football-Data.co.uk ซึ่งเป็นราคา ณ ช่วงเวลาที่บันทึก ในสถานการณ์จริงราคาอัตราต่อรองมีความผันผวนอยู่ตลอดเวลา (Odds Movement / Slippage)
    \item โครงงานไม่ได้รวมปัจจัยเชิงคุณภาพที่ไม่ปรากฏในสถิติ เช่น ข่าวสารการบาดเจ็บนาทีสุดท้าย สภาพอากาศ สภาพสนาม หรือการเปลี่ยนแปลงผู้จัดการทีม
    \item สถิติเชิงลึกรายผู้เล่นจาก Understat ครอบคลุมเฉพาะนัดในลีก ทำให้ค่าดัชนีความเหนื่อยล้ารายบุคคลอาจคลาดเคลื่อนเล็กน้อยสำหรับทีมที่เล่นฟุตบอลถ้วยหลายรายการ (ได้รับการชดเชยบางส่วนด้วยตัวแปรตารางแข่งขันระดับทีมจาก Transfermarkt)
    \item ชุดข้อมูล Transfermarkt ที่ใช้หยุดอัปเดตในเดือนกรกฎาคม 2026 นัดของฤดูกาล 2026--2027 จึงไม่มีตัวแปรตารางแข่งขันทุกรายการและแผนการเล่น และต้องเติมค่าว่างแทน
    \item ข้อมูลที่รวบรวมจากต่างแหล่งต้องใช้ตารางจับคู่ชื่อสโมสรและนักเตะมาตรฐาน (Alias Mapping) ซึ่งอาจมีค่าว่างในกรณีที่ไม่สามารถจับคู่ข้อมูลระหว่างแหล่งข้อมูลได้
    \item ผลลัพธ์และกลยุทธ์ที่ได้จัดทำขึ้นเพื่อวัตถุประสงค์ทางวิชาการและการศึกษาเชิงสถิติเท่านั้น ไม่ได้มีเจตนาในการส่งเสริมการพนันที่ไม่ถูกต้องตามกฎหมาย
\end{enumerate}

% -----------------------------------------------------------
\section{ประโยชน์ที่คาดว่าจะได้รับ}
% -----------------------------------------------------------

\begin{enumerate}[label=\arabic*)]
    \item ได้กระบวนการวิทยาการข้อมูลแบบครบวงจรสำหรับการวิเคราะห์ข้อมูลกีฬาและสถิติการเดิมพันอย่างเป็นระบบ
    \item ได้ระบบประเมินความคุ้มค่าของการเดิมพัน (Value Betting System) ที่สามารถเปรียบเทียบความน่าจะเป็นจริงของแบบจำลองกับราคาตลาด และแนะนำขนาดเงินวางเดิมพันได้อย่างเหมาะสม
    \item ได้ระบบทดสอบย้อนหลัง (Backtesting System) ที่สามารถประเมินผลกำไรและความเสี่ยงของการใช้กลยุทธ์เดิมพันรูปแบบต่าง ๆ บนข้อมูลจริงย้อนหลังฤดูกาล 2014--2015 ถึง 2025--2026
    \item ได้องค์ความรู้ในการประยุกต์ใช้เทคนิควิทยาการข้อมูล การเรียนรู้ของเครื่อง และกลยุทธ์ทางการเงิน/การวางเงินทุน เพื่อนำไปปรับใช้กับการวิเคราะห์ตลาดการแข่งขันอื่น ๆ ในอนาคต
\end{enumerate}